% Manuscript-ready insertion. Replace \ref{fig:carla-closed-loop} and
% \ref{tab:carla-*} only if the final manuscript uses different numbering.

\subsection{Closed-loop validation under converging occluded overtakes}
\label{sec:carla-validation}

\paragraph{Experimental design.}
The mechanism-level simulations in Figs.~6 and~8 were complemented by a
closed-loop CARLA validation in Town06.  The ego vehicle approached a large
rigid fire-engine surrogate while preparing to pass it from the left lane.  A
second vehicle simultaneously passed the occluder from the right lane, and
both vehicles requested a merge into the centre lane ahead of the occluder.
Nine surrounding vehicles followed an intelligent-driver-model controller to
provide matched upstream and downstream traffic.  We evaluated a
\emph{true-threat} condition containing the hidden vehicle and an
\emph{empty-shadow} condition in which the same occluded volume was empty.
DREAM and IDEAM were evaluated in a $2\times2$ matched design over five fixed
traffic seeds (20 runs in total).  The four arms of each seed shared the same
initial traffic state, reference trace, and scenario manifest.  Camera data
were disabled in the timing runs and used only in a separately replayed,
outcome-blind representative visualization.

An episode was admitted to analysis only when the complete projected
footprint of the hidden vehicle was blocked by the occluder during the
pre-specified $0.25$--$0.75$~s occlusion interval, the hidden vehicle was absent
from the ego semantic-LiDAR return before reveal, and the stored scenario and
traffic-trace hashes matched across all four arms.  The hidden vehicle was
made available to DREAM only after its first semantic-LiDAR detection; the
driver-view camera was never an input to the planner.  Each episode lasted
6~s, with CARLA physics advanced at 20~Hz, low-level actuation at 10~Hz, and
high-level optimization executed asynchronously.  Because only five matched
scene blocks were used in this initial validation, the intervals reported
below are descriptive 95\% scene-block bootstrap intervals (10,000 resamples,
fixed seed 260716), and no population-level significance claim is made.

\paragraph{Occlusion-conflict performance.}
Figure~\ref{fig:carla-closed-loop}a shows the causal sequence in the
representative scene: the entire hidden-vehicle footprint is initially
occluded, it subsequently enters the semantic-LiDAR field of view, and the
two planned merges create a conflict ahead of the occluder.  Across the five
matched true-threat scenes, DREAM produced no collision or near collision,
whereas IDEAM produced two collisions and five near collisions
(Table~\ref{tab:carla-safety-tradeoff}).  The paired differences in collision
and near-collision incidence were $-0.40$ [$-0.80$, $0.00$] and $-1.00$
[$-1.00$, $-1.00$], respectively.  Continuous margins showed the same
direction: DREAM increased the minimum oriented-box clearance to the hidden
vehicle by 1.50~m [0.92, 2.03] and the minimum two-dimensional time to
collision by 0.96~s [0.54, 1.29] relative to IDEAM.  These results provide a
closed-loop consistency check for the risk-field behavior illustrated in the
Python simulations; they are not treated as a substitute for a larger
scenario-coverage study.

\paragraph{Empty-shadow trade-off and following traffic.}
The empty-shadow arm isolates defensive behavior when no hidden vehicle is
present.  DREAM's mean ego speed was 30.09~m\,s$^{-1}$, compared with
31.39~m\,s$^{-1}$ for IDEAM, yielding an empty-shadow conservatism tax
$CT_v=1.30$~m\,s$^{-1}$ [1.22, 1.39], or approximately
4.68~km\,h$^{-1}$.  Thus, the cost of uncertainty-aware behavior was
measurable and is reported explicitly rather than characterized as negligible.
Over the 6-s observation horizon, however, the paired additional maximum
speed loss of the nearest following vehicle was only
$1.9\times10^{-6}$~m\,s$^{-1}$
[$0.4\times10^{-6}$, $3.4\times10^{-6}$], and the number of hard-braking
traffic actors was unchanged.  The additional integrated speed deficit over
all traffic actors was 0.219 vehicle-m [0.175, 0.268].  These observations
indicate no additional short-horizon follower disturbance attributable to
DREAM in the tested matched scenes, but the 6-s horizon is insufficient to
exclude longer-range traffic-wave amplification or a phantom-jam effect.

\paragraph{Asynchronous execution.}
The high-level optimizer remained slower than the nominal 100-ms planning
interval: its mean, mean per-run 95th percentile, and maximum observed
wall-clock times were 451.4, 533.1, and 579.3~ms, respectively.  At the mean
true-threat ego speed of 30.03~m\,s$^{-1}$, these durations correspond to
approximately 13.6, 16.0, and 17.4~m of travel.  The implemented scheduler
therefore coalesced pending requests and delivered 2.30 completed high-level
updates per second, dropping 76.8\% of superseded requests.  In parallel, the
low-level controller maintained 10~Hz with mean, 95th-percentile, and maximum
execution times of 1.456, 1.826, and 3.556~ms and no deadline misses; the
20-Hz physics/control loop likewise had no 50-ms deadline miss.  The low-level
layer tracked the most recent time-stamped plan and retained lane-hold and
visible-actor emergency-braking fallbacks.  The mean reveal-to-application
delay of a hidden-aware DREAM plan was 0.89~s [0.82, 0.97].  Accordingly, this
experiment demonstrates an implemented asynchronous proof of concept, but it
does not establish 10-Hz real-time feasibility of the current high-level
optimizer or a formal safety guarantee during an aggressive reveal.

\begin{table*}[t]
\centering
\caption{Matched closed-loop CARLA pilot ($n=5$ scene blocks). Intervals are
descriptive 95\% scene-block bootstrap intervals. For safety metrics, the
paired contrast is DREAM minus IDEAM; for the empty-shadow speed row,
$CT_v$ is defined as IDEAM minus DREAM.}
\label{tab:carla-safety-tradeoff}
\small
\begin{tabular}{llccc}
\toprule
Condition & Metric & DREAM & IDEAM & Paired contrast [95\% interval] \\
\midrule
True threat & Collision incidence & 0/5 & 2/5 & $-0.40$ [$-0.80$, $0.00$] \\
 & Near-collision incidence & 0/5 & 5/5 & $-1.00$ [$-1.00$, $-1.00$] \\
 & Minimum hidden-vehicle clearance (m) & 1.72 & 0.22 & $+1.50$ [0.92, 2.03] \\
 & Minimum hidden-vehicle 2-D TTC (s) & 1.20 & 0.25 & $+0.96$ [0.54, 1.29] \\
\midrule
Empty shadow & Mean ego speed (m\,s$^{-1}$) & 30.09 & 31.39 & $CT_v=1.30$ [1.22, 1.39] \\
 & Maximum nearest-follower speed loss (m\,s$^{-1}$) & 2.013816 & 2.013814 & $+1.9\times10^{-6}$ [$0.4$, $3.4$]$\times10^{-6}$ \\
 & Integrated traffic speed deficit (vehicle-m) & 62.69 & 62.47 & $+0.219$ [0.175, 0.268] \\
 & Hard-braking traffic actors (count) & 0.20 & 0.20 & 0.00 [0.00, 0.00] \\
\bottomrule
\end{tabular}
\end{table*}

\begin{table}[t]
\centering
\caption{Measured execution times for DREAM in the true-threat CARLA arm
($n=5$). The high-level row reports planner service time; the other rows
report execution time within their scheduled control loops.}
\label{tab:carla-runtime}
\small
\begin{tabular}{lrrrr}
\toprule
Layer & Period & Mean & P95 & Maximum \\
\midrule
High-level optimizer & 100 ms request & 451.4 ms & 533.1 ms & 579.3 ms \\
Low-level controller & 100 ms & 1.456 ms & 1.826 ms & 3.556 ms \\
Physics/control loop & 50 ms & 12.296 ms & 15.253 ms & 30.977 ms \\
\bottomrule
\end{tabular}
\end{table}

\begin{figure*}[t]
\centering
\includegraphics[width=\textwidth]{fig_carla_closed_loop_validation.pdf}
\caption{Closed-loop validation of the converging occluded-overtake scenario.
(a) Outcome-blind representative DREAM replay showing full initial occlusion,
semantic-LiDAR reveal, and the subsequent merge response; driver and bird's-eye
views are stacked without overlap. (b) Matched true-threat minimum clearance
to the hidden vehicle, with collision counts reported separately. (c)
Empty-shadow ego and nearest-follower velocity profiles. Quantitative results
come from camera-free timing runs; the rendered replay is illustrative.}
\label{fig:carla-closed-loop}
\end{figure*}

\begin{figure}[t]
\centering
\includegraphics[width=\columnwidth]{fig_carla_async_runtime.pdf}
\caption{Measured asynchronous execution in the DREAM true-threat arm. The
high-level optimization time exceeds its 100-ms request period, whereas the
10-Hz low-level controller and 20-Hz physics/control loop complete within
their respective deadlines. Points denote the five matched scene seeds and
bars summarize the corresponding means.}
\label{fig:carla-runtime}
\end{figure}
